% Chapter 1

\chapter{Specifica dei requisiti} % Main chapter title

\label{Chapter1} % For referencing the chapter elsewhere, use \ref{Chapter1} 

%----------------------------------------------------------------------------------------

% Define some commands to keep the formatting separated from the content 
% \newcommand{\keyword}[1]{\textbf{#1}}
% \newcommand{\tabhead}[1]{\textbf{#1}}
% \newcommand{\code}[1]{\texttt{#1}}
% \newcommand{\file}[1]{\texttt{\bfseries#1}}
% \newcommand{\option}[1]{\texttt{\itshape#1}}

%----------------------------------------------------------------------------------------

% =========================================================== %
%                 Section: General Information                %
% =========================================================== %
\section{Scopo del documento dei requisiti}
\label{chap1:sec1:general_information}
Il documento \textit{SRS}, riportato nei primi tre capitoli, si propone di accogliere in forma organica una descrizione di tutte le specifiche tecniche e funzionali che caratterizzano il sistema per la gestione di issue in progetti software. \\
Esso è stato redatto seguendo le indicazioni dello standard IEEE/ANSI 830-1998 (IEEE 1998) per i documenti dei requisiti.

% =========================================================== %
%      Section: General Requirements of Graduate School       %
% =========================================================== %
\section{Scopo del prodotto}

Il prodotto software che si vuole realizzare prende il nome di BugBoard26, una piattaforma che consentirà a un team di sviluppo di segnalare problemi relativi a un progetto, monitorarne lo stato, assegnarli a membri del team e tenere traccia delle attività di risoluzione. Il sistema deve fornire le seguenti funzionalità:

\begin{enumerate}
  \item Il sistema deve fornire un meccanismo di autenticazione semplice e sicuro, basato su email e
        password. Le informazioni gestite dall’applicazione sono critiche per l’azienda, ed è fondamentale preservarne l’integrità e la segretezza. Il sistema viene fornito con un account da amministratore già attivo, con credenziali di default. Un amministratore può creare ulteriori utenze, specificando una email, una password, e indicando se quell’utenza sarà "normale" oppure "di amministrazione".
  \item Tutti gli utenti autenticati possono segnalare una issue indicando almeno un titolo e una
        descrizione. Alcuni utenti potrebbero voler specificare anche una priorità e sarebbe gradita la possibilità di allegare un’immagine. Le issue possono essere di diverso tipo: question (per richieste di chiarimenti), bug (per segnalare malfunzionamenti), documentation (per segnalare problemi relativi alla documentazione), e feature (per indicare la richiesta o il suggerimento di nuove funzionalità). Le issue create sono inizialmente nello stato "todo".
  \item Il sistema deve offrire una vista riepilogativa delle issue, con la possibilità di filtrare        o ordinare i risultati in base a criteri come tipologia, stato, priorità o altri parametri         rilevanti.
  \item Ogni bug dovrebbe poter essere assegnato (da un amministratore) a un membro del team. Quando
        viene assegnato un bug a un utente, quest’ultimo riceve una notifica. Gli utenti possono
        visualizzare i bug loro assegnati.
  \item È necessaria una funzione di ricerca che consenta di trovare bug in base a parole chiave.
  \item Il sistema dovrebbe suggerire automaticamente a chi assegnare un bug, basandosi sul carico di
        lavoro corrente degli utenti.
  \item Un amministratore può contrassegnare un bug come duplicato di un altro. Questo comporta la chiusura immediata del bug duplicato.

\end{enumerate}

\section{Riferimenti}
Qui di seguito sono elencati i documenti e gli indirizzi web a cui fa riferimento la presente specifica dei requisiti, nonché i 
link al codice sorgente e al repository Docker
\begin{itemize}
    \item Slide del corso di Ingegneria del software 2025/2026 fornite dai professori P.Tramontana e B.Breve.
    \item Ian, Sommerville. \textit{Ingegneria del software}. 8\textsuperscript{a} edizione, Pearson Paravia Bruno Mondadori S.p.A., 2007.
    \item Modello \href{https://ieeexplore.ieee.org/document/720574}{IEEE/ANSI 830-1998} per strutturare il documento dei requisiti.
    \item Enciclopedia online \href{https://it.wikipedia.org/wiki/Pagina_principale}{Wikipedia}. 
    \item \href{https://github.com/ubiestubi/ProgettoIngSoft}{Repository del back-end}: la guida
      all'installazione è nel file \texttt{README.md}.
    \item \href{https://github.com/GiandomenicoIameo/BugBoard26/}{Repository del front-end}: la guida
      all'installazione è nel file \texttt{README.md} e i pacchetti eseguibili sono nella sezione Release.
\end{itemize}






